\documentclass[10pt,a4paper]{amsart}
\usepackage[margin=19mm]{geometry}
\usepackage{amsmath,amssymb,mathtools}
\usepackage[scr=rsfs,cal=euler]{mathalfa}
\usepackage{xfrac}
\usepackage[unicode,hidelinks,bookmarks=false]{hyperref}
\newcommand{\ie}{i.e.\ }
\newcommand{\cf}{cf.\ }
\newcommand{\resp}{respectively\ }
\newcommand{\Subsection}{\subsection}
\providecommand{\nobreakdash}{\nobreak\hbox{-}}
\setlength{\emergencystretch}{2em}
\multlinegap0pt
\usepackage{amssymb}
\newtheorem{lemma}{Lemma}
\title{A note on the self-intersection of rational unicuspidal curve with one Puiseux pair}
\author{Weimin Chen, Yusupjan Ouyang, Sam Silver}
\date{Source: arXiv:2608.06534v1 (CC BY 4.0)}
\begin{document}
\maketitle

\noindent\emph{Source and licence.} A note on the self-intersection of rational unicuspidal curve with one Puiseux pair. Version arXiv:2608.06534v1 (CC BY 4.0), distributed by arXiv under the Creative Commons Attribution 4.0 International licence, http://creativecommons.org/licenses/by/4.0/, as recorded in the arXiv licence metadata for this exact version and shown on the abstract page https://arxiv.org/abs/2608.06534v1. Full licence notice: \texttt{licenses/arxiv-2608.06534-CC-BY-4.0.txt}.\par\smallskip

\noindent\textit{Editorial scope.} Counted unit: the article's lemma lemma (source0.tex lines 343--347). The article's complete original proof of it is delivered with it (source0.tex lines 349--393, one \texttt{proof} environment of 45 lines). No mathematics is written, repaired or re-derived: the statement and the proof are the article's own lines, in the article's own notation, and no retained line is substituted or re-flowed.  No further source-local context is needed: every cross-reference of every retained line resolves inside the two delivered spans.  Nothing was omitted from any retained span, so the package carries no omission record.  No retained line contains a citation, so the wrapper carries no bibliography and no bibliography entry.  No retained line contains a figure environment, an image-inclusion command or drawing code, so no image is delivered and the package declares no asset dependency.  Editorial formatting only, in addition: an AMS \texttt{amsart} preamble that restates the article's own declarations and macros used by the retained lines, namely the environments \texttt{lemma} on separate counters, together with the standard packages those lines need.  The retained environments print in this document as Lemma 1 in order of appearance, whereas the article prints the same items with the numbering of its own full document.  The complete native source of the article as distributed in the arXiv e-print for this exact version is bundled unchanged as \texttt{sources/207006/source0.tex}.
\begin{lemma}
Assume $1<p<q$ where $gcd(p,q)=1$, and set $\Delta_{p,q}:=B_{p,q}-R_{p,q}$. 
Then $\Delta_{p,q}\in \{0,1\}$. Moreover, $\Delta_{p,q}=1$
if and only if either $q=p+1$ or $p>2$ and $q=-1\pmod{p}$. 
\end{lemma}

\begin{proof}
Case (1): suppose $q=p+1$. By a direct calculation, one has $B_{p,q}=B_{p,p+1}=p+1$ and 
$R_{p,q}=R_{p,p+1}=p$, which implies that $\Delta_{p,q}=1$.

\vspace{2mm}

Case (2): suppose $q\neq p+1$, and we write $q=kp+r$ where $1\leq r<p$. 

\vspace{2mm}

(i): assume $r=1$, where $k>1$ is necessarily true. In this case, by Lemma 2.2, then Lemma  2.1, one has $B_{p,q}=B_{p, p+1}-1=R_{p,p+1}=R_{p,q}$, which implies that $\Delta_{p,q}=0$ in this case.

\vspace{2mm} 

(ii): assume $1<r=p-1$, where $p>2$ is necessarily true. Note that in this case, $p>2$ and $q=-1\pmod{p}$ is true, so we need to prove $\Delta_{p,q}=1$. To see this, we observe that with $r=p-1$, the condition $r>2$ and $p=-1\pmod{r}$ cannot be satisfied. Hence by Lemma 2.2, then
Lemma 2.1, one has 
$$
B_{p,q}=B_{r,p}=B_{p-1, p}=R_{p-1,p}+1=R_{r,p}+1=R_{p,q}+1,
$$ 
which implies $\Delta_{p,q}=1$. (Here we used the fact that $\Delta_{p-1,p}=1$ from Case (1).)

\vspace{2mm} 

(iii): assume $1<r<p-1$. We need to show $\Delta_{p,q}=0$. 
By induction, we assume the lemma is true for $(r,p)$. With this understood, 
we observe that by Lemma 2.2,
\[
  B_{p,q}=
  \begin{cases}
    B_{r,p}-1, & \text{if } r>2 \text{ and } p= -1\pmod r,\\
    B_{r,p}, & \text{if otherwise, i.e., either } r=2 \text{ or } p\neq -1\pmod r.
  \end{cases}
\]
In the former case, since the lemma is true for $(r,p)$, we have $B_{r,p}=R_{r,p}+1=R_{p,q}+1$,
as $\Delta_{r,p}=1$ by the induction assumption and $R_{r,p}=R_{p,q}$ by Lemma 2.1. Consequently, 
$$
\Delta_{p,q}=B_{p,q}-R_{p,q}=B_{r,p}-1-R_{p,q}=0.
$$
Likewise, in the latter case, we have $B_{r,p}=R_{r,p}=R_{p,q}$ instead, as in this case, since
$p\neq r+1$ is also true, one has $\Delta_{r,p}=0$ by the induction assumption. It follows that 
$$
\Delta_{p,q}=B_{p,q}-R_{p,q}=B_{r,p}-R_{p,q}=0
$$
as well. This finishes the proof. 
\end{proof}

\end{document}
